% Source printed pp. 446-459 / PDF pp. 454-467.
\section{半稳定向量丛}\label{src:9.1}
\begin{epigraph}
纯数学最主要的应用是使你快乐.
\epigraphauthor{Hendrik Lenstra}
\end{epigraph}
我们介绍光滑曲线上向量丛的半稳定性这一基本概念. 乍看之下, 这个定义似乎缺乏动机, 但我们将说明它具有若干理想的性质:
\begin{itemize}
\item 每个向量丛都有唯一的滤过---Harder--Narasimhan 滤过---其因子半稳定 (定理~\bookref{9.1.21}); 每个半稳定向量丛都有一个 (非典范的) 滤过---Jordan--H\"older 滤过---其因子稳定 (命题~\bookref{9.1.32}).
\item 半稳定性是开放且有界的条件 (定理~\bookref{9.2.4} 和 \bookref{9.2.8}).
\item 半稳定向量丛的代数叠 $\cBun_{r,d}^{\mathrm{ss}}$ 普遍闭 (定理~\bookref{9.3.1}).
\item 存在射影良模空间 $\cBun_{r,d}^{\mathrm{ss}}\to M_{r,d}^{\mathrm{ps}}$, 它将 S-等价向量丛等同, 并诱导多重稳定向量丛 (对应于 $\cBun_{r,d}^{\mathrm{ss}}$ 的闭点) 与 $M_{r,d}^{\mathrm{ps}}(\Bbbk)$ 之间的双射 (定理~\hyperref[thm:B]{B} 和命题~\bookref{9.1.35}).
\item 对每个 $r>0$ 和 $d$, 都存在半稳定向量丛 (定理~\bookref{9.1.36}). 事实上, 也存在稳定向量丛 (定理~\bookref{9.2.16}).
\end{itemize}
在 \S\bookref{9.1.6-subsection}, 我们以椭圆曲线上的向量丛为例进行研究.
\subsection{光滑曲线上凝聚层的回顾}\label{src:9.1.1-subsection}
设 $C$ 是任意域 $\Bbbk$ 上的光滑几何连通射影曲线. $C$ 上的任何凝聚层 $F$ 都有挠子层 $F_{\mathrm{tors}}$, 它在开子集 $U\subseteq C$ 上的截面是 $F(U)$ 中挠元素组成的子集 $F(U)_{\mathrm{tors}}$. 商 $F/F_{\mathrm{tors}}$ 是向量丛, 且由于
\[
\operatorname{Ext}_{\mathcal O_C}^1(F/F_{\mathrm{tors}},F_{\mathrm{tors}})
=\mathrm H^1(C,F_{\mathrm{tors}}\otimes(F/F_{\mathrm{tors}})^\vee)=0
\]
(因为 $\dim(\operatorname{Supp}F_{\mathrm{tors}})=0$), 短正合列 $0\to F_{\mathrm{tors}}\to F\to F/F_{\mathrm{tors}}\to0$ 分裂, 从而 $F\cong F_{\mathrm{tors}}\oplus F/F_{\mathrm{tors}}$.

若 $F$ 是秩为 $\operatorname{rk}F=r$ 的向量丛, 则定义它的行列式为线丛 $\det F:=\bigwedge^rF$, 次数为 $\deg F:=\deg(\det F)$. 对凝聚层 $F\cong F_{\mathrm{tors}}\oplus F/F_{\mathrm{tors}}$, 定义 $F$ 的次数为 $\deg F:=\mathrm h^0(C,F_{\mathrm{tors}})+\deg(F/F_{\mathrm{tors}})$. 定义凝聚层 $F$ 的秩为 $\operatorname{rk}F:=\dim_{\kappa(\xi)}F\otimes\kappa(\xi)$, 其中 $\xi\in C$ 为一般点. 挠层的秩为 $0$. 若 $0\to F'\to F\to F''\to0$ 是短正合列, 则 $\deg F=\deg F'+\deg F''$ 且 $\operatorname{rk}F=\operatorname{rk}F'+\operatorname{rk}F''$. 换言之, 次数与秩延拓为从 $C$ 的 Grothendieck 群出发的群同态
\[
\deg,\operatorname{rk}:K(C)\to\mathbb Z
\]
\cite[习题 II.6.10]{Har77}.
\begin{exercise}[简单]\label{src:9.1.1}
若 $F_1$ 和 $F_2$ 是 $C$ 上的向量丛, 证明 $\deg(F_1\otimes F_2)=\deg F_1\operatorname{rk}F_2+\operatorname{rk}F_1\deg F_2$.
\end{exercise}
线丛的 Riemann--Roch 定理 (\bookref{6.1.2}) 可推广到凝聚层.
\begin{theorem}[凝聚层的 Riemann--Roch 定理]\label{src:9.1.2}
设 $C$ 是代数闭域 $\Bbbk$ 上的光滑连通射影曲线. 若 $F$ 是 $C$ 上的凝聚层, 则
\[
\chi(C,F)=\deg(F)+\operatorname{rk}(F)(1-g).
\]
\end{theorem}
\begin{proof}
利用 Euler 示性数, 秩与次数对短正合列均可加, 对秩作归纳. 若 $F$ 是挠层, 则由次数的定义, $\chi(C,F)=\deg F$. 若 $\operatorname{rk}(F)>0$, 则存在一个线丛子层 $L\subseteq F$, 满足 $\operatorname{rk}(F/L)=\operatorname{rk}(F)-1$. 由归纳假设, 结论对 $F/L$ 成立; 由 Riemann--Roch 定理 (\bookref{6.1.2}), 结论对 $L$ 成立, 因而对 $F$ 也成立.
\end{proof}
\begin{exercise}[简单]\label{src:9.1.3}
证明 $C$ 上秩为 $r$, 次数为 $d$ 的向量丛 $F$ 的 Hilbert 多项式为 $P(t):=rt+(d+r(1-g))$.
\end{exercise}
\begin{exercise}[简单]\label{src:9.1.4}
若 $F_1,F_2$ 是 $C$ 上次数为 $d_i$, 秩为 $r_i$ 的向量丛, 证明
\[
\dim\operatorname{Ext}_{\mathcal O_C}^1(F_1,F_2)-\dim\operatorname{Ext}_{\mathcal O_C}^0(F_1,F_2)
=r_1r_2(g-1)+(d_1r_2-d_2r_1).
\]
特别地, 若 $F$ 是秩为 $r$, 次数为 $d$ 的向量丛, 则 $\dim\operatorname{Ext}_{\mathcal O_C}^1(F,F)-\dim\operatorname{Ext}_{\mathcal O_C}^0(F,F)=r^2(g-1)$.
\end{exercise}
我们还回顾 Serre 对偶 (\bookref{6.1.3}): 对 $C$ 上凝聚层 $F$, 存在自然同构
\[
\mathrm H^1(C,F)\cong\Hom_{\mathcal O_C}(F,\Omega_C)^\vee.
\]
\textbf{例: $\mathbb P^1$ 上的向量丛.} 一个经典定理断言, $\mathbb P^1$ 上的向量丛分裂为线丛的直和.
\begin{theorem}[Birkhoff--Grothendieck 定理]\label{src:9.1.5}
若 $E$ 是任意域 $\Bbbk$ 上 $\mathbb P^1$ 上的向量丛, 则存在唯一的整数 $d_1\geq\cdots\geq d_r$, 使得 $E\cong\bigoplus_{i=1}^r\mathcal O_{\mathbb P^1}(d_i)$.
\end{theorem}
\begin{proof}
对秩 $r$ 作归纳. 若 $r=1$, 则每个线丛都同构于 $\mathcal O(d)$ \cite[命题 II.6.4]{Har77}. 若 $E$ 是秩 $r>0$ 的向量丛, 令 $d_1$ 为满足 $\mathrm H^0(\mathbb P^1,E(-d_1))\ne0$ 的最大整数. $E(-d_1)$ 的非零截面定义单射 $\mathcal O(d_1)\to E$ 及短正合列
\begin{equation}\label{src:9.1.6}
0\to\mathcal O(d_1)\to E\to E'\to0.
\end{equation}
我们断言 $E'$ 是向量丛: 否则, 核 $K:=\ker(E\to E'/E'_{\mathrm{tors}})$ 是满足 $\mathcal O(d_1)\subsetneq K$ 的线丛, 从而 $K\cong\mathcal O(d)$, 其中 $d>d_1$, 与 $d_1$ 的选取矛盾. 由归纳假设, 可写为 $E'=\bigoplus_{i=2}^r\mathcal O(d_i)$, 其中 $d_2\geq\cdots\geq d_r$. 扭曲上述列得到短正合列 $0\to\mathcal O(-1)\to E(-d_1-1)\to\bigoplus_{i=2}^r\mathcal O(d_i-d_1-1)\to0$. 由于 $\mathrm H^0(\mathbb P^1,E(-d_1-1))=0$ 且 $\mathrm H^0(\mathbb P^1,\mathcal O(-1))=\mathrm H^1(\mathbb P^1,\mathcal O(-1))=0$, 对所有 $i\geq2$, 有 $\mathrm H^0(\mathbb P^1,\mathcal O(d_i-d_1-1))=0$. 因而 $d_i-d_1-1<0$, 等价地, $d_i\leq d_1$. 为证明 \bookeqref{9.1.6} 分裂, 应用 $\Hom_{\mathcal O_{\mathbb P^1}}(-,\mathcal O(d_1))$ 并计算得
\[
\operatorname{Ext}_{\mathcal O_{\mathbb P^1}}^1(E',\mathcal O(d_1))
=\bigoplus_{i\geq2}\mathrm H^1(\mathbb P^1,\mathcal O(d_1-d_i))=0,
\]
因为 $d_1-d_i\geq0$. 唯一性留给读者. 另见 \cite[引理 4.4.1]{LP97} 和 \cite[定理 1.3.1]{HL10}.

另一种证明使用矩阵分解, 其思路与 Birkhoff 的论证 \cite{Bir09} 相似. 因为 $\mathbb A^1$ 上的每个向量丛都平凡, $\mathbb P^1$ 上的向量丛 $E$ 可由过渡函数 $g\in\mathrm{GL}_r(\Bbbk[x]_x)$ 指定, 它在 $E|_{\mathbb P^1\setminus0}$ 和 $E|_{\mathbb P^1\setminus\infty}$ 的自同构作用下唯一; 即存在双射
\[
\{\mathbb P^1\text{ 上的向量丛}\}/{\sim}\longleftrightarrow
\mathrm{GL}_r(\Bbbk[x])\backslash\mathrm{GL}_r(\Bbbk[x]_x)/\mathrm{GL}_r(\Bbbk[1/x]).
\]
只需证明: 对任意 $g\in\mathrm{GL}_r(\Bbbk[x]_x)$, 存在 $A\in\mathrm{GL}_r(\Bbbk[x])$ 和 $B\in\mathrm{GL}_r(\Bbbk[1/x])$, 使得 $AgB=\operatorname{diag}(x^{d_1},\ldots,x^{d_r})$, 其中整数 $d_1\geq\cdots\geq d_r$. 这一线性代数事实留作习题.
\end{proof}
\begin{remark}\label{src:9.1.7}
Grothendieck 对约化群 $G$ 在 $\mathbb P^1$ 上的 $G$-丛证明了类似结论 \cite{Gro57a}, 在 $G=\mathrm{GL}_r$ 时得到上述结论.
\end{remark}
\subsection{半稳定性的定义与初步性质}\label{src:9.1.2-subsection}
\begin{definition}[斜率]\label{src:9.1.8}
任意域 $\Bbbk$ 上光滑几何连通射影曲线 $C$ 上的凝聚层 $F$ 的\newterm{斜率}{slope} 定义为
\[
\mu(F):=\frac{\deg F}{\operatorname{rk}F}\in\mathbb Q\cup\{\infty\},
\]
约定当 $F$ 为挠层时 $\mu(F)=\infty$.
\end{definition}
\begin{remark}[斜率的可视化]\label{src:9.1.9}
可以将凝聚层 $F$ 在复平面中表示为点 $Z(F):=-\deg(F)+\operatorname{rk}(F)i$.\footnote{读者或许觉得将它表示为点 $(\operatorname{rk}F,\deg F)$ 更自然. 将荷写为 $Z(F)=-\deg(F)+\operatorname{rk}(F)i\in\mathbb C$, 是受 Bridgeland 稳定性的启发: $(\mathsf{Coh}(C),Z)$ 定义 $D^b(C)$ 上的一个稳定性条件.} 按此约定, $E$ 的斜率大于 $F$ 的斜率, 当且仅当 $Z(E)$ 的辐角大于 $Z(F)$ 的辐角. 函数 $Z(F)$ 有时称为\newterm{荷}{charge}, 对短正合列可加. 这一观点可用来直观理解基本性质, 例如, 对子层 $E\subseteq F$, $\mu(E)\leq\mu(F)$ 当且仅当 $\mu(F/E)\geq\mu(F)$.
\end{remark}
\bookfiguretarget{9.1.10}
\begin{figure}[H]
\centering
\begin{tikzpicture}[scale=0.85,>=Stealth]
\draw (-2.8,0)--(3.2,0) node[right] {$-\deg$};
\draw (0,-0.25)--(0,3.5) node[above] {$\operatorname{rk}$};
\draw[->] (0,0)--(0.65,1.75) node[above left] {$Z(E)$};
\draw[->] (0,0)--(2.8,2.05) node[right] {$Z(F)$};
\draw[blue,dashed,->] (0.65,1.75)--(2.8,2.05) node[midway,above] {$Z(F/E)$};
\end{tikzpicture}
\bookfigurehref{9.1.10}
\caption{斜率的可视化: 对子层 $E\subseteq F$, 有 $Z(F)=Z(E)+Z(F/E)$.}\label{src:9.1.10}
\end{figure}
\begin{definition}[半稳定性/稳定性/多重稳定性]\label{src:9.1.11}
域 $\Bbbk$ 上光滑几何连通射影曲线 $C$ 上的非零凝聚层 $F$ 称为
\begin{enumerate}[(1)]
\item \newterm{半稳定}{semistable}, 若对每个非零子层 $0\ne E\subseteq F$, 有 $\mu(E)\leq\mu(F)$;
\item \newterm{稳定}{stable}, 若对每个非零真子层 $0\ne E\subsetneq F$, 有 $\mu(E)<\mu(F)$;
\item \newterm{多重稳定}{polystable}, 若 $F\cong\bigoplus_iF_i$, 其中每个 $F_i$ 稳定且 $\mu(F_i)=\mu(F)$.
\end{enumerate}
若 $F$ 不半稳定, 则满足 $\mu(E)>\mu(F)$ 的子层 $E\subseteq F$ 称为\newterm{不稳定化子层}{destabilizing subsheaf}. 若 $F$ 半稳定但不稳定, 则称 $F$ \newterm{严格半稳定}{strictly semistable}.
\end{definition}
\begin{example}[\texorpdfstring{$C=\mathbb P^1$}{C=P1}]\label{src:9.1.12}
$\mathbb P^1$ 上的稳定向量丛只有线丛. 向量丛 $E\cong\bigoplus_i\mathcal O(d_i)$ 半稳定, 当且仅当它多重稳定, 当且仅当所有 $d_i$ 相等.
\end{example}
\begin{lemma}\label{src:9.1.13}
\begin{enumerate}[(1)]
\item 若 $F$ 是秩为 $r$, 次数为 $d$ 的半稳定向量丛, 且 $\gcd(r,d)=1$, 则 $F$ 稳定.
\item 半稳定凝聚层是向量丛.
\item 向量丛 $F$ 半稳定 (分别地, 稳定), 当且仅当对每个非零真子丛 $F'\subsetneq F$ (即 $F'$ 是子层且 $F/F'$ 为向量丛), 有 $\mu(F')\leq\mu(F)$ (分别地, $<$).
\item 对 $C$ 上的线丛 $L$, 向量丛 $F$ 半稳定当且仅当 $F\otimes L$ 半稳定.
\end{enumerate}
\end{lemma}
\begin{proof}
对 (1), 若子丛 $0\ne E\subseteq F$ 满足 $\mu(E)=\mu(F)$, 则存在整数 $k\geq1$, 使得 $\operatorname{rk}F=k\operatorname{rk}E$ 且 $\deg F=k\deg E$. 由于 $\gcd(r,d)=1$, 有 $k=1$ 且 $F=E$. 对 (2), 若 $F$ 不是向量丛, 则挠子层 $F_{\mathrm{tors}}\subseteq F$ 是不稳定化子层. 对 (3), 若 $E\subseteq F$ 是不稳定化子层, 商为 $Q=F/E$, 则 $E':=\ker(F\to Q\to Q/Q_{\mathrm{tors}})$ 是包含 $E$ 的 $F$ 的子丛, 且 $\operatorname{rk}(E')=\operatorname{rk}(E)$, $\deg(E')\geq\deg(E)$. 因为 $\mu(E')\geq\mu(E)$, $E'\subseteq F$ 是不稳定化子丛. 对 (4), 注意 $\deg(F\otimes L)=\deg(F)+\operatorname{rk}(F)\cdot\deg(L)$ 且 $\mu(F\otimes L)=\mu(F)+\deg(L)$. 所以, 对任何子层 $E\subseteq F$, 有
\[
\mu(F)-\mu(E)=\mu(F\otimes L)-\mu(E\otimes L),
\]
这蕴含 $F$ 半稳定当且仅当 $F\otimes L$ 半稳定.
\end{proof}
\begin{exercise}[简单]\label{src:9.1.14}
\begin{enumerate}[(a)]
\item 证明凝聚层 $F$ 半稳定, 当且仅当对每个凝聚商层 $F\twoheadrightarrow Q$, 有 $\mu(Q)\geq\mu(F)$.
\item 证明向量丛 $F$ 半稳定当且仅当 $F^\vee$ 半稳定.
\end{enumerate}
\end{exercise}
\begin{lemma}[简单]\label{src:9.1.15}
\begin{enumerate}[(1)]
\item 若 $E,F$ 是半稳定向量丛, 且 $\mu(E)>\mu(F)$, 则 $\Hom_{\mathcal O_C}(E,F)=0$.
\item 若 $E,F$ 是稳定向量丛, 且 $\mu(E)=\mu(F)$, 则非零同态 $E\to F$ 是同构.
\end{enumerate}
\end{lemma}
\begin{proof}
对 (1), 设 $E\to F$ 是非零映射, 以 $K$ 和 $I=E/K$ 分别表示其核与像. 因为 $E$ 半稳定, $\mu(K)\leq\mu(E)$. 所以
\[
\begin{aligned}
\deg(I)=\mu(I)\operatorname{rk}(I)&=\mu(E)\operatorname{rk}(E)-\mu(K)\operatorname{rk}(K)\\
&\geq\mu(E)\operatorname{rk}(E)-\mu(E)\operatorname{rk}(K)\\
&=\mu(E)\operatorname{rk}(I),
\end{aligned}
\]
从而 $\mu(I)\geq\mu(E)>\mu(F)$. 除非 $I=0$, 这与 $F$ 的半稳定性矛盾. 对 (2), 若 $E\to F$ 是斜率 $\alpha=\mu(E)=\mu(F)$ 的稳定层之间的非零映射, 则核 $K$ 的斜率也为 $\alpha$. 因为 $E$ 稳定且 $K\subsetneq E$, $K=0$. 类似地, 像 $I$ 的斜率为 $\alpha$, 且由于 $F$ 稳定, $I=F$. 因此 $E\to F$ 是同构.
\end{proof}
\begin{corollary}\label{src:9.1.16}
设 $\Bbbk$ 代数闭. 若 $F$ 是稳定向量丛, 则 $\underline{\operatorname{Aut}}(F)\cong\mathbb G_m$. 更一般地, 若 $F=\bigoplus F_i^{\oplus r_i}$ 是多重稳定向量丛, 每个 $F_i$ 稳定且斜率为 $\mu(F)$, 并且 $i\ne j$ 时 $F_i\not\cong F_j$, 则 $\underline{\operatorname{Aut}}(F)\cong\prod_i\mathrm{GL}_{r_i}$.
\end{corollary}
\begin{proof}
若 $F$ 是稳定丛, 由 (2), $\operatorname{Aut}(F)$ 是除代数. 因为 $\Bbbk$ 代数闭, $\operatorname{Aut}(F)\cong\Bbbk^*$. 群概形的同构 $\underline{\operatorname{Aut}}(F)\cong\mathbb G_m$ 由一个简单的形式论证得到 (见命题~\bookref{9.2.20}(2)). 对最后一个断言, 因为 $i\ne j$ 时 $\Hom_{\mathcal O_C}(F_i,F_j)=0$, 有 $\operatorname{Aut}(F)\cong\prod_i\operatorname{Aut}(F_i^{\oplus r_i})=\prod_i\mathrm{GL}_{r_i}$.
\end{proof}
\begin{exercise}[有益练习, 简单]\label{src:9.1.17}
证明固定斜率的半稳定向量丛组成的范畴是对扩张封闭的 Abel 范畴.
\end{exercise}
多重稳定向量丛的自同构群约化 (推论~\bookref{9.1.16}), 但并非每个半稳定向量丛都有这一性质.
\begin{exercise}[中等]\label{src:9.1.18}
证明严格半稳定向量丛的自同构群可能不约化.
\end{exercise}
\begin{remark}[Narasimhan--Seshadri 定理]\label{src:9.1.19}
设 $C$ 是 $\mathbb C$ 上的光滑连通射影曲线, 将它视为 Riemann 曲面. 设 $\widetilde C\to C$ 为万有覆盖. 若 $\pi_1(C)\to\mathrm{GL}_n(\mathbb C)$ 是表示, 则 $\pi_1(C)$ 对角地作用于 $\widetilde C\times\mathbb C^n$, 商 $(\widetilde C\times\mathbb C^n)/\pi_1(C)$ 是 $C$ 上某向量丛的全空间. Narasimhan 与 Seshadri 证明此对应诱导 $\pi_1(C)$ 的不可约酉表示与 $C$ 上稳定向量丛之间的双射 \cite{NS65}.
\end{remark}
\begin{remark}[张量积下的半稳定性]\label{src:9.1.20}
Narasimhan 与 Seshadri 利用稳定性与不可约酉表示之间的对应, 证明在 $\mathbb C$ 上两个半稳定向量丛的张量积半稳定. 这一性质在正特征下不成立.
\end{remark}
